% Generated by uscribe -- do not edit by hand
\documentclass[letterpaper,12pt]{article}
\usepackage[T1]{fontenc}
\usepackage[utf8]{inputenc}
\usepackage{graphicx}
\usepackage{hyperref}
\usepackage{xcolor}
\usepackage{fancyvrb}
\usepackage{booktabs}
\usepackage{longtable}
\IfFileExists{multirow.sty}{\usepackage{multirow}}{}
\hypersetup{colorlinks=true,linkcolor=blue,urlcolor=blue}
\begin{document}
\begin{titlepage}
\null\vfil\vskip 40pt
\begin{center}
{\LARGE\bfseries Raw Sockets and Packet Layouts in Unicon \par}
\vskip 3em
{\large\bfseries Jafar Al-Gharaibeh \par}
\vskip 1.0em
{\bfseries Unicon Technical Report: 29\par}
\vskip 1.0em
{\large\bfseries July 2026\par}
\vskip 5.5em
{\bfseries Abstract}\par
\begin{quote}
Unicon already opened TCP and UDP sockets through open(), but IP-level protocols such as ICMP and IGMP were not supported. This report describes raw sockets: mode "nr" opens SOCK\_RAW, proto= selects the IP protocol number, and writes()/receive() carry upper-layer bytes on the same datagram path UDP already uses. The runtime does not assemble protocol payloads. Package net supplies PacketSpec, a declarative engine for fixed binary layouts, plus ICMP echo and IGMPv2/v3 classes and session wrappers.
\end{quote}
\begin{quote}
{\bfseries Keywords:} Unicon, raw sockets, SOCK\_RAW, ICMP, IGMP, PacketSpec, runtime, technical report.
\end{quote}
\vskip 0.8in
{\large\bfseries
Unicon Project\\
\url{http://unicon.org}\\
\ \
University of Idaho\\
Department of Computer Science\\
Moscow, ID, 83844, USA
}
\end{center}
\vfil\null
\end{titlepage}
\setcounter{page}{1}
\section{1. Introduction}
\label{utr29-sec-intro}

This report describes raw IP sockets in Unicon: how a handle is opened, how datagrams are sent and received, how protocol bytes are built and decoded, and what remains out of scope. It is formatted as a Unicon Technical Report [1]. Socket attributes (\texttt{ttl}, \texttt{iface}, \texttt{join}, and the rest) are the multicast facility [11]; this report adds mode \texttt{nr} and the packet library on top of that. The language-facing reference lives in \emph{Programming with Unicon} [8] (Chapter 5, "Raw Sockets" and "PacketSpec and package net"; Chapter 14, \texttt{uping} and \texttt{uigmp}) and the language reference (\texttt{open} mode \texttt{nr}, attributes \texttt{proto} and \texttt{hdrincl}). Automated plumbing tests live in \texttt{tests/posix/raw\_sock.icn}.


Raw sockets are part of the POSIX network facility, not an optional library bind. A build without POSIX networking has no \texttt{n} modes at all. Opening \texttt{SOCK\_RAW} usually needs elevated privileges (root or administrator); that is an operating-system restriction, not a Unicon one.

\section{2. Motivation}
\label{utr29-sec-motivation}

TCP and UDP cover most application protocols. Diagnostic and control traffic sits one layer down: ICMP echo [12], IGMP membership [13] [14], OSPF, GRE, PIM. Those messages are IP protocol payloads, not port-addressed transport streams. ICMP, IGMP, and other raw IP protocols were not supported.


The rest of Unicon networking already treats \texttt{open()} as the constructor and \texttt{Attrib()} as the option verb. Raw sockets follow that pattern instead of adding \texttt{rawsocket()} or a parallel packet API. Payload construction is a library problem: the runtime delivers bytes, and \texttt{PacketSpec} names the fields.

\section{3. Design principles}
\label{utr29-sec-principles}

\textbf{\texttt{r} after \texttt{n} is the socket type, like \texttt{u} for UDP.} Mode letter \texttt{r} remains ordinary read mode unless a preceding \texttt{n} has already marked the handle as a network socket. Then \texttt{"nr"} means \texttt{SOCK\_RAW}, the same way \texttt{"nu"} means \texttt{SOCK\_DGRAM}. \texttt{"rn"} is not a raw socket: \texttt{r} is consumed as read mode before \texttt{n} is seen.


\textbf{\texttt{proto=} is required and is fixed at \texttt{socket()} time.} The third argument to \texttt{socket(2)} is the IP protocol number. There is no useful default: ICMP, IGMP, and OSPF are different sockets. A missing, empty, or unknown \texttt{proto=} is error 1310 (\texttt{bad~socket~attribute}). \texttt{Attrib(f,~"proto")} cannot query or change it afterwards.


\textbf{The runtime does not build protocol payloads.} Writes send whatever bytes the caller supplies. Receives return whatever the kernel delivered, often including an IPv4 header on Linux. \texttt{import~net} and \texttt{PacketSpec} (or hand-assembled strings) fill that gap.


\textbf{Datagram I/O matches UDP.} Raw sockets are not connected. \texttt{writes()} uses \texttt{sendto(2)} with the destination saved at \texttt{open()}. Incoming packets use \texttt{receive()}, which returns a record with \texttt{addr} and \texttt{msg}, not \texttt{read()} / \texttt{reads()}.


\textbf{Two layers for callers.} Session wrappers (\texttt{Ping}, \texttt{Igmp}) hide mode letters for the common cases. PacketSpec classes are there when the program must craft or classify bytes itself.

\section{4. Opening a raw socket}
\label{utr29-sec-open}
\begin{Verbatim}[fontsize=\small,frame=single,commandchars=\\\{\}]
f := open("8.8.8.8", "nr", "proto=icmp", "ttl=64") |
     stop(&errortext)
\end{Verbatim}


The first argument is a bare host or address, not \texttt{host:port}. Raw destinations are hosts. A single \texttt{host:port} form is accepted when there is exactly one colon and a non-empty port, but ICMP and IGMP do not use ports. IPv6 literals contain colons; they are treated as hosts and are not split as \texttt{host:port}.


\texttt{proto=} takes a symbolic name or an integer 0--255:

\begin{longtable}{ll}
\hline
Name & Protocol \\
\hline
\endfirsthead
\hline
Name & Protocol \\
\hline
\endhead
\hline
\endfoot
\hline
\endlastfoot
\texttt{icmp} & \texttt{IPPROTO\_ICMP} \\
\texttt{icmpv6} / \texttt{icmp6} & \texttt{IPPROTO\_ICMPV6} (fails if the host has no constant) \\
\texttt{igmp} & \texttt{IPPROTO\_IGMP} \\
\texttt{tcp} / \texttt{udp} & \texttt{IPPROTO\_TCP} / \texttt{IPPROTO\_UDP} \\
\texttt{gre} & 47 \\
\texttt{ospf} & 89 \\
\texttt{pim} & \texttt{IPPROTO\_PIM} \\
\texttt{raw} & \texttt{IPPROTO\_RAW} \\
\texttt{0} .. \texttt{255} & that protocol number \\
\end{longtable}


Names are case-insensitive. Unknown names and numbers outside 0--255 fail at \texttt{open()} with 1310.


If the kernel refuses \texttt{SOCK\_RAW}, \texttt{open()} fails and \texttt{\&errortext} is the system message (typically "Operation not permitted" or "Permission denied"). That is success of the Unicon plumbing and a privilege problem on the host.

\begin{Verbatim}[fontsize=\small,frame=single,commandchars=\\\{\}]
procedure main()
   local f
   if f := open("8.8.8.8", "nr") then
      write("oops missing proto")
   else
      write("missing-proto: ", &errortext)
   if f := open("8.8.8.8", "nr", "proto=nope") then
      write("oops bad proto")
   else
      write("bad-proto: ", &errortext)
   if f := open("8.8.8.8", "nr", "proto=icmp") then
      write("raw-icmp: ok")
   else
      write("raw-icmp: ", &errortext)
end
\end{Verbatim}


Output:

\begin{Verbatim}[fontsize=\small,frame=single,commandchars=\\\{\}]
missing-proto: bad socket attribute
bad-proto: bad socket attribute
raw-icmp: Operation not permitted
\end{Verbatim}


The last line is the unprivileged case. With privileges it prints \texttt{raw-icmp:~ok} (or the test suite's \texttt{plumbing-ok}).

\section{5. Attributes}
\label{utr29-sec-attrs}

Trailing \texttt{open()} arguments are the same \texttt{name=value} socket attributes used for TCP and UDP [11]. Booleans are exactly \texttt{yes} or \texttt{no}. \texttt{ttl} and \texttt{iface} were named for this sharing: hop limit and interface are not multicast-only knobs.

\begin{longtable}{ll}
\hline
Attribute & Role on a raw socket \\
\hline
\endfirsthead
\hline
Attribute & Role on a raw socket \\
\hline
\endhead
\hline
\endfoot
\hline
\endlastfoot
\texttt{proto=} & Required. Consumed at \texttt{socket()}. Not queryable. \\
\texttt{hdrincl=yes} & Sets \texttt{IP\_HDRINCL} so writes may include a complete IPv4 \\
\texttt{ttl=} & Unicast and multicast hop limits (\texttt{IP\_TTL} / \\
\texttt{rcvbuf=} / \texttt{sndbuf=} & Socket buffer sizes in bytes. \\
\texttt{broadcast=} / \texttt{mcastloop=} & Same meaning as UDP. \\
\texttt{join=} / \texttt{leave=} / \texttt{source=} / \texttt{iface=} & Multicast membership, including SSM \texttt{source@group} \\
\end{longtable}


\texttt{Attrib(f,~"ttl=4")} and \texttt{Attrib(f,~"ttl")} use the same names after \texttt{open()}. \texttt{proto} remains create-time only. \texttt{hdrincl} can be read back as \texttt{yes} / \texttt{no}.

\section{6. Send and receive}
\label{utr29-sec-io}
\begin{Verbatim}[fontsize=\small,frame=single,commandchars=\\\{\}]
writes(f, icmp_bytes)
r := receive(f)
write(r.addr, "  ", *r.msg, " bytes")
\end{Verbatim}


\texttt{writes()} calls \texttt{sock\_write()}, which for \texttt{SOCK\_RAW} (and UDP) does \texttt{sendto} to the \texttt{addrinfo} saved at \texttt{open()}. There is no \texttt{connect(2)}. Changing the destination means opening another handle (or, with \texttt{hdrincl=yes}, putting the destination in a supplied IP header).


\texttt{receive()} peeks with \texttt{MSG\_PEEK} into a 2K buffer, then consumes with \texttt{recvfrom}. If the peek fills 2K, the consume uses a 64K buffer so a large datagram is not truncated. The result is a record: \texttt{addr} is the peer (often \texttt{a.b.c.d:0} for raw IPv4), \texttt{msg} is the byte string.


On many IPv4 stacks a raw receive buffer starts with the IP header, then the protocol message. Linux ICMP is the usual example. Strip it before decoding ICMP or IGMP (\hyperref[utr29-sec-packetspec]{section 7}). IPv6 raw sockets typically deliver the upper layer only; there is no IPv6 header-strip helper yet.


\texttt{select()} works on the handle. Session wrappers use it to implement timeouts around \texttt{receive()}.


As with UDP, \texttt{read()} / \texttt{reads()} are the wrong verbs for datagrams. Use \texttt{receive()}.

\section{7. PacketSpec}
\label{utr29-sec-packetspec}

Module \texttt{packetspec} (package \texttt{net}, \texttt{uni/lib/packetspec.icn}) is a small declarative engine for fixed binary layouts. Protocol classes inherit \texttt{PacketSpec} and override \texttt{define()}.


Layer-1 helpers pack and unpack network-order integers (\texttt{u8} / \texttt{u16} / \texttt{u32} / \texttt{u64} and the matching \texttt{get\_u*} readers) and IPv4 addresses (\texttt{ipv4} / \texttt{to\_ipv4} / \texttt{get\_ipv4}). \texttt{inet\_cksum()} is RFC 1071 [15].

\subsection{7.1 Declaring a layout}
\label{utr29-sec-ps-fields}
\begin{Verbatim}[fontsize=\small,frame=single,commandchars=\\\{\}]
class IcmpEchoRequest : IcmpPacket()
   method define()
      field("type", 1, 8)
      field("code", 1, 0)
      checksum_field("checksum")
      field("id", 2)
      field("seq", 2)
   end
end
\end{Verbatim}


\texttt{field(name,~size,~defval)} is a network-order integer. \texttt{bytes\_field} is opaque fixed-width bytes. \texttt{ipv4\_field} accepts a dotted-quad, integer, or four raw bytes on build and returns a dotted-quad on decode. \texttt{checksum\_field} is a 16-bit placeholder filled at \texttt{build()} time over the header plus payload. Pass \texttt{exclude\_off} / \texttt{exclude\_len} (1-based) to omit a range from the sum --- OSPF Authentication and PIM Register are the documented cases; those layouts are not in the library yet.


\texttt{require(name)} marks a field that \texttt{build()} must see via \texttt{set()} or the optional extra table. \texttt{relax()} / \texttt{build(...,~lax)} skip that check for incomplete test packets. \texttt{strict()} and \texttt{clear()} turn checking back on.

\subsection{7.2 Build, decode, describe}
\label{utr29-sec-ps-build}

Preferred call shape (no string-keyed tables required). This needs no raw socket:

\begin{Verbatim}[fontsize=\small,frame=single,commandchars=\\\{\}]
import net

procedure main()
   local pkt, v
   write(IcmpEchoRequest().describe())
   write(IcmpEchoRequest().set("id", 1).set("seq", 2).describe())
   pkt := IcmpEchoRequest().set("id", 1).set("seq", 2).build(, "hi")
   write("len=", *pkt)
   v := IcmpEchoRequest().decode(pkt)
   write("type=", v["type"], " id=", v["id"], " seq=", v["seq"])
end
\end{Verbatim}


Output:

\begin{Verbatim}[fontsize=\small,frame=single,commandchars=\\\{\}]
type:int=8  code:int=0  checksum:checksum(auto)  id:int=0  seq:int=0
type:int=8  code:int=0  checksum:checksum(auto)  id:int=1  seq:int=2
len=10
type=8 id=1 seq=2
\end{Verbatim}


\texttt{set()} and \texttt{set\_payload()} are chainable. \texttt{build(extra,~payload,~lax)} overlays an optional table, appends payload, and returns the complete byte string with checksum filled. \texttt{decode(s)} returns a table of field names to values. \texttt{describe()} is a one-line summary for the REPL: names, kinds, defaults, required markers, and current \texttt{set()} values.

\subsection{7.3 IPv4 receive helpers}
\label{utr29-sec-ps-ip}

\texttt{strip\_ipv4(pkt)} uses the header IHL to return the upper-layer message. \texttt{ipv4\_ttl(pkt)} and \texttt{ipv4\_src(pkt)} read TTL and source from the same leading header. ICMP and IGMP classes call \texttt{strip\_ipv4} before matching.

\section{8. ICMP echo}
\label{utr29-sec-icmp}

\texttt{uni/lib/icmp.icn} adds echo request (type 8) and echo reply (type 0) [12] on top of \texttt{PacketSpec}:

\begin{Verbatim}[fontsize=\small,frame=single,commandchars=\\\{\}]
PacketSpec
  IcmpPacket          strip_ip / match
    IcmpEchoRequest   type 8; echo(id, seq, payload)
    IcmpEchoReply     type 0; match_echo(pkt, id, seq)
Ping                  session wrapper
\end{Verbatim}


\texttt{Ping(host)} opens \texttt{proto=icmp} with TTL 64 and a 2000 ms default wait. \texttt{echo()} sends a request (56 zero bytes unless a payload is given), waits with \texttt{select()}, and succeeds with a decode table plus \texttt{rtt} (milliseconds, real), \texttt{rtt\_us}, \texttt{ttl} (from the IPv4 header), and \texttt{addr}. Round-trip time uses \texttt{gettimeofday()}, not \texttt{\&time}, because CPU time barely advances while blocked in \texttt{select()}.

\begin{Verbatim}[fontsize=\small,frame=single,commandchars=\\\{\}]
import net

procedure main()
   local p, v
   p := Ping("8.8.8.8") | stop(&errortext)
   write(p.describe())
   if v := p.echo() then
      write("rtt=", v["rtt"], " from ", v["addr"])
   p.close()
end
\end{Verbatim}


\texttt{set\_timeout(ms)} changes the default wait. An explicit sequence number and payload can be passed to \texttt{echo()}; otherwise the session increments \texttt{seq}.


The same exchange with an explicit socket:

\begin{Verbatim}[fontsize=\small,frame=single,commandchars=\\\{\}]
import net

procedure main()
   local f, req, reply, id, pkt, r
   req := IcmpEchoRequest()
   reply := IcmpEchoReply()
   id := iand(?100000 | 1, 16rFFFF)
   f := open("8.8.8.8", "nr", "proto=icmp", "ttl=64") |
        stop(&errortext)
   pkt := req.set("id", id).set("seq", 1).build(, "unicon")
   writes(f, pkt)
   r := receive(f)
   if reply.match_echo(r.msg, id, 1) then
      write("reply from ", r.addr)
   close(f)
end
\end{Verbatim}


Other ICMP types (destination unreachable, time exceeded, and so on) are not in the library. \texttt{proto=icmpv6} opens an ICMPv6 socket, but there are no Neighbor Discovery or echo layouts yet.

\section{9. IGMP}
\label{utr29-sec-igmp}

\texttt{uni/lib/igmp.icn} covers IGMPv2 (RFC 2236) and IGMPv3 (RFC 3376):

\begin{Verbatim}[fontsize=\small,frame=single,commandchars=\\\{\}]
PacketSpec
  IgmpPacket
    Igmpv2Query / Igmpv2Report / Igmpv2Leave
    Igmpv3Query / Igmpv3Report
Igmp                     session wrapper
\end{Verbatim}


IGMPv2 messages are an eight-octet header. Report and leave \texttt{require()} a group. IGMPv3 queries add flags, QQIC, and a source list. IGMPv3 reports carry group records: ASM join/leave and SSM source filters (\texttt{join\_sources}, \texttt{allow}, \texttt{block}).


\texttt{Igmp()} opens \texttt{proto=igmp} bound at \texttt{0.0.0.0}. \texttt{recv(ms)} waits, then returns a table with \texttt{kind}, \texttt{addr}, \texttt{msg}, \texttt{len}, and decoded fields when recognized. \texttt{kind} is one of \texttt{v3\_report}, \texttt{v3\_query}, \texttt{v2\_query}, \texttt{v2\_report}, \texttt{v2\_leave}, or \texttt{unknown}.


PacketSpec classes work without a raw socket:

\begin{Verbatim}[fontsize=\small,frame=single,commandchars=\\\{\}]
import net

procedure main()
   local pkt, v
   write(Igmpv2Report().describe())
   write(Igmpv3Report().describe())
   pkt := Igmpv2Report().report("239.1.1.1")
   write("v2 report len=", *pkt)
   v := Igmpv2Report().decode(pkt)
   write("type=", v["type"], " group=", v["group"])
   pkt := Igmpv3Report().join("239.1.1.1")
   write("v3 join len=", *pkt)
end
\end{Verbatim}


Output:

\begin{Verbatim}[fontsize=\small,frame=single,commandchars=\\\{\}]
type:int=22  max_resp_time:int=0  checksum:checksum(auto)  group:ipv4(required)
type:int=34  reserved1:int=0  checksum:checksum(auto)  reserved2:int=0  nrecords:int=0
v2 report len=8
type=22 group=239.1.1.1
v3 join len=16
\end{Verbatim}


Opening a live IGMP socket needs the same privileges as ICMP:

\begin{Verbatim}[fontsize=\small,frame=single,commandchars=\\\{\}]
import net

procedure main()
   local g, v
   g := Igmp() | stop(&errortext)
   write(g.describe())
   if v := g.recv(1000) then
      write(v["kind"], " from ", v["addr"])
   g.report("239.1.1.1")    # IGMPv2 membership report
   g.join("239.1.1.1")      # IGMPv3 ASM join
   g.close()
end
\end{Verbatim}


Session send helpers: \texttt{report} / \texttt{leave} / \texttt{query} (v2), \texttt{join} / \texttt{leave3} (v3 ASM), and \texttt{send(bytes)} for a hand-built packet. SSM source-filter joins use the PacketSpec class, not the session wrapper:

\begin{Verbatim}[fontsize=\small,frame=single,commandchars=\\\{\}]
pkt := Igmpv3Report().join_sources("239.1.1.1", ["192.0.2.1"])
g.send(pkt)
\end{Verbatim}


Exported constants (\texttt{IGMP\_MEMBERSHIP\_QUERY}, \texttt{IGMP\_V3\_MEMBERSHIP\_REPORT}, \texttt{IGMP\_CHANGE\_TO\_EXCLUDE\_MODE}, and the rest) are initialized on first \texttt{IgmpPacket} construction.

\section{10. Sample programs}
\label{utr29-sec-progs}

\texttt{uni/progs/uping} and \texttt{uni/progs/uigmp} are teaching clients, not replacements for \texttt{ping(8)} or \texttt{tcpdump(1)}. They need \texttt{SOCK\_RAW} privileges:

\begin{Verbatim}[fontsize=\small,frame=single,commandchars=\\\{\}]
unicon uping
unicon uigmp
sudo ./uping host [count]
sudo ./uigmp [seconds]
sudo ./uping -v host    # also show open()/Ping API demos
sudo ./uigmp -v         # also show open()/Igmp API demos
\end{Verbatim}


Default mode uses the session wrappers. \texttt{-v} / \texttt{-d} also runs the lower-level \texttt{open(...,~"nr",~...)} plus PacketSpec path and prints \texttt{describe()} output.

\begin{Verbatim}[fontsize=\small,frame=single,commandchars=\\\{\}]
sudo ./uping 192.0.2.1
\end{Verbatim}


Output:

\begin{Verbatim}[fontsize=\small,frame=single,commandchars=\\\{\}]
PING 192.0.2.1 (192.0.2.1): 56 data bytes
64 bytes from 192.0.2.1: icmp_seq=0 ttl=118 time=18.390 ms
64 bytes from 192.0.2.1: icmp_seq=1 ttl=118 time=16.245 ms
64 bytes from 192.0.2.1: icmp_seq=2 ttl=118 time=15.533 ms
64 bytes from 192.0.2.1: icmp_seq=3 ttl=118 time=17.311 ms
--- 192.0.2.1 ping statistics ---
4 packets transmitted, 4 packets received, 0.0% packet loss
round-trip min/avg/max = 15.533/16.869/18.390 ms
\end{Verbatim}

\begin{Verbatim}[fontsize=\small,frame=single,commandchars=\\\{\}]
sudo ./uigmp
\end{Verbatim}


Output:

\begin{Verbatim}[fontsize=\small,frame=single,commandchars=\\\{\}]
listening on IGMP

36 bytes from 192.0.2.10: IGMPv3 query group=0.0.0.0 max_resp=100 qrv=2
36 bytes from 192.0.2.10: IGMPv3 query group=239.1.1.1 max_resp=100 qrv=2
36 bytes from 192.0.2.20: IGMPv3 report \{239.1.1.1 CHANGE_TO_EXCLUDE\}
\end{Verbatim}

\section{11. Runtime implementation}
\label{utr29-sec-impl}

The work is in the existing POSIX socket code, not a new file type.


\textbf{Mode letter.} \texttt{src/runtime/fsys.r}: after \texttt{n} has set \texttt{Fs\_Socket}, \texttt{r} / \texttt{R} sets \texttt{sock\_type~=~SOCK\_T\_RAW} and does not set the ordinary read bit. Without \texttt{Fs\_Socket}, \texttt{r} remains \texttt{Fs\_Read}. \texttt{e} (crypto raw-key material) is a later check on the same letter; \texttt{"nr"} is a socket, \texttt{"er"} is a crypto handle.


\textbf{Create.} \texttt{src/runtime/rposix.r}: \texttt{sock\_connect} (and the listen/bind path) require \texttt{proto=} before \texttt{socket()}. \texttt{uni\_getaddrinfo} allows a missing port for \texttt{SOCK\_T\_RAW}. The real protocol number is passed as the third argument to \texttt{socket()}; \texttt{getaddrinfo} is told protocol 0 so it does not override it.


\textbf{I/O.} UDP and raw share \texttt{saddrs[]}: the \texttt{addrinfo} from \texttt{open()} is kept until \texttt{sock\_close()}. \texttt{sock\_write} \texttt{sendto}s there. \texttt{sock\_recv} accepts \texttt{SOCK\_DGRAM} and \texttt{SOCK\_RAW}. \texttt{proto=} is skipped in \texttt{setsockopt} (already consumed). \texttt{hdrincl=} sets \texttt{IP\_HDRINCL}.


\textbf{Errors.} Unknown or missing socket attributes use error 1310 with the offending string (\texttt{proto}, \texttt{proto=nope}, \texttt{hdrincl=maybe}). Kernel failures use the system \texttt{\&errortext}.


There is no \texttt{\&features} flag named "raw sockets". If POSIX sockets exist, mode \texttt{nr} exists.

\section{12. Status and remaining work}
\label{utr29-sec-status}

The language surface described here is implemented: \texttt{"nr"}, required \texttt{proto=}, \texttt{hdrincl=}, datagram \texttt{writes} / \texttt{receive}, PacketSpec, ICMP echo, IGMPv2/v3, \texttt{uping}, and \texttt{uigmp}.


\textbf{More PacketSpec layouts.} GRE, OSPF, and PIM names are accepted at \texttt{socket()} and the checksum helper documents their exclude ranges, but there are no protocol classes yet.


\textbf{ICMPv6.} \texttt{proto=icmpv6} opens the socket. Echo, Neighbor Discovery, and an IPv6 header strip are not in \texttt{net}.


\textbf{IPv6 receive metadata.} \texttt{strip\_ipv4} / \texttt{ipv4\_ttl} assume a leading IPv4 header. IPv6 raw receives need their own helpers.


\textbf{\texttt{read()} on raw handles.} UDP clears the read bit so \texttt{read()} fails. Raw currently keeps \texttt{Fs\_Read} because it shares the TCP-like status assignment. The documented API is still \texttt{receive()}. Aligning the status bits with UDP would make misuse fail earlier.


\textbf{Privileged tests.} \texttt{tests/posix/raw\_sock.icn} checks plumbing without sending packets. A live ICMP echo test would need root in CI and a reachable responder; it is not in the suite.


\textbf{Windows.} Opening \texttt{SOCK\_RAW} is possible with administrator rights, but Windows historically restricts which protocols a raw socket may use. Treat \texttt{uping} / \texttt{uigmp} as UNIX-first tools.


\textbf{Header inclusion.} \texttt{hdrincl=yes} is wired. There is no PacketSpec IPv4 header class to go with it; callers who craft full datagrams build those bytes themselves.

\section{Appendix: test coverage}
\label{utr29-sec-coverage}

\texttt{tests/posix/raw\_sock.icn} does not require \texttt{SOCK\_RAW} success. It checks that:

\begin{itemize}
\item 
missing, empty, and unknown \texttt{proto=} fail with 1310
\item 
\texttt{proto=icmp}, \texttt{gre}, \texttt{ospf}, and numeric \texttt{89} reach \texttt{socket()} (success or a privilege denial both print \texttt{plumbing-ok})
\item 
\texttt{hdrincl=maybe} is rejected
\item 
without a preceding \texttt{n}, mode \texttt{"r"} still opens a file for reading
\end{itemize}

Expected output is \texttt{tests/posix/stand/raw\_sock.std}:

\begin{Verbatim}[fontsize=\small,frame=single,commandchars=\\\{\}]
missing-proto: bad socket attribute
bad-proto: bad socket attribute
empty-proto: bad socket attribute
raw-icmp: plumbing-ok
raw-gre: plumbing-ok
raw-ospf: plumbing-ok
raw-numeric: plumbing-ok
bad-hdrincl: bad socket attribute
read-mode-r: ok
\end{Verbatim}


On an unprivileged account the \texttt{raw-*} lines still print \texttt{plumbing-ok}: the test treats both a successful \texttt{socket()} and a privilege denial as passing plumbing. The posix suite picks up every \texttt{*.icn} in that directory.

\section{References}
\label{utr29-references}
\begin{thebibliography}{99}
\bibitem[1]{cite-jeffery-utr15}
Clinton Jeffery. "How to Write a Unicon Technical Report". doc/utr/utr15.tex. 2013.

\bibitem[2]{cite-libssh}
The libssh Project. "libssh -- The SSH Library". https://www.libssh.org/. 2026.

\bibitem[3]{cite-jeffery-utr13}
Clinton Jeffery. "The Unicon Messaging Facilities". doc/utr/utr13.odt.

\bibitem[4]{cite-openssl}
The OpenSSL Project. "OpenSSL -- Cryptography and SSL/TLS Toolkit". https://www.openssl.org/. 2026.

\bibitem[5]{cite-algharaibeh-utr26}
Jafar Al-Gharaibeh. "Native SSH and SFTP Support in Unicon". doc/utr/utr26.rst. 2026.

\bibitem[6]{cite-rfc2104}
H. Krawczyk and M. Bellare and R. Canetti. "HMAC: Keyed-Hashing for Message Authentication". RFC 2104. 1997.

\bibitem[7]{cite-rfc1112}
S. Deering. "Host Extensions for IP Multicasting". RFC 1112. 1989.

\bibitem[8]{cite-jeffery-pwu}
Clinton Jeffery and Shamim Mohamed and Jafar Al-Gharaibeh. "Programming with Unicon". doc/book/. 2026.

\bibitem[9]{cite-algharaibeh-utr29}
Jafar Al-Gharaibeh. "Raw Sockets and Packet Layouts in Unicon". doc/utr/utr29.rst. 2026.

\bibitem[10]{cite-rfc4607}
H. Holbrook and B. Cain. "Source-Specific Multicast for IP". RFC 4607. 2006.

\bibitem[11]{cite-algharaibeh-utr28}
Jafar Al-Gharaibeh. "Multicast and Socket Attributes in Unicon". doc/utr/utr28.rst. 2026.

\bibitem[12]{cite-rfc792}
J. Postel. "Internet Control Message Protocol". RFC 792. 1981.

\bibitem[13]{cite-rfc2236}
W. Fenner. "Internet Group Management Protocol, Version 2". RFC 2236. 1997.

\bibitem[14]{cite-rfc3376}
B. Cain and S. Deering and I. Kouvelas and B. Fenner and A. Thyagarajan. "Internet Group Management Protocol, Version 3". RFC 3376. 2002.

\bibitem[15]{cite-rfc1071}
R. Braden and D. Borman and C. Partridge. "Computing the Internet Checksum". RFC 1071. 1988.

\end{thebibliography}
\end{document}
